% Hướng dẫn sử dụng cho người lên kế hoạch sản xuất ERP Sao Việt Nhật.
% Biên dịch: xelatex huong-dan-ke-hoach-san-xuat.tex (chạy hai lần cho mục lục).
% Người đọc: docs/personas/ke-hoach-san-xuat.md. Khai báo nằm ở ../huong-dan-cau-hinh/cau-hinh-he-thong.tex,
% phần đơn hàng bán ở ../huong-dan-kinh-doanh/huong-dan-su-dung-kinh-doanh.tex.
% Ảnh chụp bằng kịch bản chạy lại được (chup-ui.mjs + c02..c07), dữ liệu dựng bằng thao tác thật trên dev.
\documentclass[11pt,a4paper]{article}
\usepackage{fontspec}
\usepackage[vietnamese]{babel}
\setmainfont{Times New Roman}
\newfontfamily\tieude{Arial}
\usepackage[a4paper,margin=1.9cm,top=1.7cm]{geometry}
\usepackage{graphicx}
\usepackage{enumitem}
\usepackage{xcolor}
\usepackage{titlesec}
\usepackage{tocloft}
\usepackage{needspace}
\usepackage{fancyhdr}
\usepackage[hidelinks,destlabel=true]{hyperref}
\usepackage[most]{tcolorbox}
\usepackage{booktabs}
\usepackage{tabularx}
\graphicspath{{hinh/}}

\definecolor{gach}{HTML}{C2410C}
\definecolor{xamnen}{HTML}{F4F4F5}
\definecolor{xamchu}{HTML}{52525B}

\setlist{nosep,leftmargin=1.5em}
\setlength{\parskip}{0.4em}
\setlength{\parindent}{0pt}
\setlength{\emergencystretch}{2em}

\titleformat{\section}{\tieude\Large}{Chương \thesection.}{0.4em}{}[{\color{black!25}\titlerule[0.4pt]}]
\titlespacing*{\section}{0pt}{1.2em}{0.6em}
\titleformat{\subsection}{\tieude\large}{\thesubsection}{0.6em}{}
\titlespacing*{\subsection}{0pt}{1.1em}{0.3em}
\renewcommand{\cftsecpresnum}{Chương }
\setlength{\cftsecnumwidth}{5.2em}
\renewcommand{\cftsecfont}{\tieude}
\renewcommand{\cftsecpagefont}{\tieude}
\setlength{\cftbeforesecskip}{0.5em}
\addto\captionsvietnamese{\renewcommand{\contentsname}{Mục lục}}
\renewcommand{\cfttoctitlefont}{\tieude\large}
\setlength{\cftbeforetoctitleskip}{0pt}
\setlength{\cftaftertoctitleskip}{0pt}

\pagestyle{fancy}
\fancyhf{}
\fancyhead[L]{\tieude\footnotesize\color{xamchu}Hướng dẫn sử dụng Kế hoạch sản xuất}
\fancyhead[R]{\tieude\footnotesize\color{xamchu}ERP Sao Việt Nhật}
\fancyfoot[C]{\tieude\footnotesize\thepage}
\renewcommand{\headrulewidth}{0.4pt}

\newtcolorbox{boicanh}{colback=xamnen,colframe=xamnen,boxrule=0pt,arc=3pt,
  left=9pt,right=9pt,top=6pt,bottom=6pt,before skip=4pt,after skip=8pt}

\definecolor{vang}{HTML}{FEF08A}
\newcommand{\oo}[1]{{\setlength{\fboxsep}{1.2pt}\colorbox{vang}{\strut #1}}}
\newcommand{\duong}[1]{{\tieude\small #1}}
\newcommand{\so}[1]{{\tieude\bfseries\small\textcolor{white}{\colorbox{red!80!black}{\kern-2pt#1\kern-2pt}}}}
\newcommand{\khung}[2]{{\setlength{\fboxsep}{0pt}\setlength{\fboxrule}{0.5pt}%
  \fcolorbox{black!20}{white}{\includegraphics[width=\dimexpr#1-1pt\relax]{#2}}}}
\newcommand{\q}[1]{\textquotedbl#1\textquotedbl{}}
\newcommand{\canh}[3]{%
  \par\noindent
  \begin{minipage}[t]{#1\linewidth}\vspace{0pt}\khung{\linewidth}{#2}\end{minipage}\hfill
  \begin{minipage}[t]{\dimexpr0.97\linewidth-#1\linewidth\relax}\vspace{0pt}\setlength{\parskip}{0.5em}#3\end{minipage}\par\medskip}
\newcommand{\rong}[2][\linewidth]{\par\noindent\begin{center}\khung{#1}{#2}\end{center}\vspace{-0.6em}}
% Ảnh cắt bớt: \cat[rộng]{trái dưới phải trên (px)}{ảnh}.
\newcommand{\cat}[3][\linewidth]{\par\noindent\begin{center}{\setlength{\fboxsep}{0pt}\setlength{\fboxrule}{0.5pt}%
  \fcolorbox{black!20}{white}{\includegraphics[width=\dimexpr#1-1pt\relax,trim=#2,clip]{#3}}}\end{center}\vspace{-0.6em}}
% Tên tài liệu khác.
\newcommand{\tl}[1]{\textit{#1}}
% Liên kết sang một mục của tài liệu Cấu hình và tài liệu Kinh doanh (đích là tên \label bên đó, nhờ destlabel).
\newcommand{\cauhinh}{../huong-dan-cau-hinh/cau-hinh-he-thong.pdf}
\newcommand{\kinhdoanh}{../huong-dan-kinh-doanh/huong-dan-su-dung-kinh-doanh.pdf}
\definecolor{lienket}{HTML}{1E40AF}
\newcommand{\kbn}[2]{\href{\cauhinh\##1}{\textcolor{lienket}{#2}}}
\newcommand{\kb}[2]{\tl{Hướng dẫn cấu hình hệ thống}, mục \kbn{#1}{#2}}
\newcommand{\kdn}[2]{\href{\kinhdoanh\##1}{\textcolor{lienket}{#2}}}
\newcommand{\kd}[2]{\tl{Hướng dẫn sử dụng phòng Kinh doanh}, mục \kdn{#1}{#2}}

\begin{document}

\thispagestyle{empty}
{\tieude\small\color{xamchu} ERP Sao Việt Nhật}\par\vspace{3pt}
{\tieude\fontsize{22}{26}\selectfont Hướng dẫn sử dụng Kế hoạch sản xuất}\par\vspace{4pt}
{\small Nhận đơn, lên lệnh, lo vật tư, xếp lịch máy, gia công ngoài, theo dõi hàng}\par\vspace{2pt}
{\small\color{xamchu} Cập nhật 05/11/2026}\par\vspace{6pt}
\noindent{\color{black!25}\rule{\linewidth}{0.4pt}}\par\vspace{4pt}

Tài liệu cho người lên kế hoạch của xưởng, dùng hằng ngày. Đi theo đúng đường một lệnh sản xuất: đơn rơi xuống,
tạo lệnh, rà lệnh cho đủ số, lo vật tư, xếp lịch máy và phát hành xuống xưởng, đưa hàng đi gia công ngoài, rồi
theo dõi tới lúc KCS đóng lệnh. Người đọc đã biết nghề in nên tài liệu không giải nghĩa tay sách, bù hao hay khổ
tờ, chỉ nói phần mềm làm gì với chúng.

Tài liệu viết theo bộ quyền của vai Kế hoạch SX, phạm vi Tất cả ở mọi màn:

\begin{itemize}
\item Kế hoạch sản xuất, Kế hoạch vật tư, Yêu cầu mua hàng: Xem và Thao tác.
\item Xếp lịch: Xem, Thao tác và Phát hành lịch.
\item Hồ sơ lệnh sản xuất, Theo dõi sản xuất: Xem.
\item Xem danh mục Công đoạn, Thiết bị \& Máy móc, Khuôn.
\end{itemize}

Thiếu ô nào thì nút tương ứng không hiện hoặc bấm bị từ chối. Quản trị bật quyền ở \kb{phan-quyen-ma-tran}{Mở ma
trận quyền} và gán vai ở mục \kbn{phan-quyen-gan-vai}{Gán vai và lập nhóm}. Chữ màu xanh là liên kết, bấm vào mở đúng
mục bên tài liệu đó.

Ảnh chụp từ phần mềm thật, dữ liệu là đơn DH003 của khách Café Katinat với ba lệnh LSX26-0007, 0008, 0009. Ô đỏ
đánh số là chỗ cần bấm hoặc cần đọc, số khớp với lời giải thích bên cạnh.

\bigskip
\tableofcontents

% ============================================================
\Needspace*{16\baselineskip}
\section{Đường đi của một lệnh}\label{duong-di}

Mỗi hàng của sơ đồ là một phòng, mỗi ô là một bước. Hàng thứ hai là việc của Kế hoạch sản xuất, các hàng còn lại
là việc phòng khác làm, kế hoạch chỉ theo dõi. Chữ trên mũi tên là điều kiện hoặc thứ được chuyển sang bước sau.

\rong[0.92\linewidth]{khsx-00-luong.png}

Đơn hàng bán do Kinh doanh chốt. Đơn cần cọc thì phải chờ Kế toán thu đủ cọc, đơn không cần cọc thì sang thẳng
(xem \kd{don-coc}{Chờ cọc}). Đơn sang là rơi vào \oo{Hàng chờ tiếp nhận} ở màn Kế hoạch sản xuất. Kế hoạch tạo lệnh,
rà cho đủ số rồi bấm \oo{Sẵn sàng lập kế hoạch}. Lệnh sẵn sàng thì lo vật tư (giữ hàng có sẵn, thiếu thì đề nghị mua),
kéo lệnh vào lịch máy và bấm \oo{Phát hành}. Phát hành xong, tổ nhận việc ở màn Việc của tổ, kho nhận yêu cầu xuất
vật tư, bước thuê ngoài chờ mang đi gia công. Tổ ghi sản lượng, KCS kiểm, gửi thành phẩm vào kho và đóng lệnh.

Lệnh giao trọn cho nhà gia công đi đường tắt: không xếp lịch, không xuống tổ, hàng về thì nhập thẳng kho thành phẩm
(chương \ref{gia-cong}).

\Needspace*{10\baselineskip}
\subsection{Trạng thái của lệnh}\label{trang-thai}

\begin{itemize}
\item Nháp: lệnh vừa tạo hoặc vừa mở lại để sửa. Sửa mọi thứ được.
\item Sẵn sàng: đủ dữ liệu, vào khay Chờ xếp lịch. Muốn sửa thì bấm Mở lại để sửa, lệnh về Nháp.
\item Đã phát hành: đã xuống xưởng. Tổ thấy việc, kho thấy yêu cầu xuất. Chuỗi công đoạn bị khoá.
\item Đã đóng: KCS đã đóng lệnh.
\end{itemize}

Kéo lệnh vào lịch không đổi trạng thái. Lệnh có giờ mà chưa phát hành vẫn là Sẵn sàng, ở Xếp lịch nó hiện chip
\q{Có lịch, chưa phát hành}. Trạng thái \q{Đã lập kế hoạch} chỉ gặp ở lệnh nằm trong bài ghép.

\Needspace*{8\baselineskip}
\subsection{Màn hình dùng hằng ngày}\label{man-hinh}

Cả năm màn nằm trong nhóm \duong{Sản xuất} của thanh bên trái:

\begin{itemize}
\item \duong{Kế hoạch sản xuất}: hàng chờ tiếp nhận, danh sách lệnh, chi tiết từng lệnh.
\item \duong{Kế hoạch vật tư}: vật tư của mọi lệnh đang chạy, giữ hàng, đề nghị mua.
\item \duong{Xếp lịch}: lịch máy, khay chờ xếp, phát hành.
\item \duong{Theo dõi sản xuất}: hàng của từng đơn tới đâu.
\item \duong{Hồ sơ lệnh sản xuất}: lệnh đã phát hành, sản lượng từng bước, in phiếu công nghệ.
\end{itemize}

Kế hoạch không làm bốn việc: đóng lệnh (KCS), tạm dừng việc (tổ bấm ở Việc của tổ), lập đơn mua (Thu mua), xuất
nhập kho (thủ kho).

% ============================================================
\Needspace*{16\baselineskip}
\section{Nhận đơn và tạo lệnh}\label{nhan-don}

Đơn chốt xong (và đủ cọc nếu cần cọc) là tự sang sản xuất. Người kế hoạch nhận thông báo ngay, tab
\oo{Hàng chờ tiếp nhận} có chấm đỏ. Mở \duong{Sản xuất > Kế hoạch sản xuất}.

\rong{khsx-10-hang-cho.png}

\so{1} Tab Hàng chờ tiếp nhận. Chấm đỏ còn đó khi còn đơn chưa xem.
\so{2} Mỗi dòng là một đơn còn mặt hàng chưa có lệnh. Cột Chuyển lúc là giờ đơn sang sản xuất. Ngày giao đỏ là đã
quá hạn, vàng là còn ba ngày trở xuống. Cột Lưu ý sản xuất chép từ đơn.
\so{3} Tiến độ lệnh: bao nhiêu dòng của đơn đã có lệnh. Đơn có đủ lệnh thì rời hàng chờ.

\Needspace*{14\baselineskip}
\subsection{Ngăn Lệnh dự kiến}\label{lenh-du-kien}

Bấm một dòng, ngăn bên phải mở ra với mỗi mặt hàng của đơn là một lệnh dự kiến.

\rong[0.9\linewidth]{khsx-11-lenh-du-kien.png}

\so{1} Số trong ngăn lấy sẵn từ bài tính giá của mặt hàng. Đây là đề xuất, sau khi tạo lệnh vẫn sửa được.
\so{2} Tick dòng muốn tạo lệnh. Dòng đã có lệnh thì hiện mã lệnh và không tick được nữa.
\so{3} Vào máy, Bù hao, Giấy nguyên, Con/tờ, Kẽm, Lượt in: số của bài in.
\so{4} Công đoạn: chuỗi bước từ bài tính giá. Bước có hình xe tải là bước thuê ngoài.
\so{5} Bấm \oo{Xác nhận tạo N lệnh sản xuất}.

Mấy trường hợp cần để ý trong ngăn:

\begin{itemize}
\item Số lượng trên phiếu tính giá khác số lượng đơn: ngăn báo, lệnh lấy số lượng đơn.
\item Dòng \q{kho đang có N}: kho thành phẩm đang có sẵn N cái của mặt hàng này, cân nhắc trước khi làm thêm.
\item Dòng \q{Chưa có bài tính giá}: mặt hàng chưa gắn phiếu tính giá nên không có số để bung. Hỏi Kinh doanh.
\item Sản phẩm nhiều phần (bìa và ruột, hộp và lót): mỗi phần là một lệnh, gom chung một nhóm trong ngăn.
\end{itemize}

Tạo xong, màn chuyển sang tab \oo{Lệnh SX}, lọc sẵn theo đơn vừa làm, đầu bảng có băng \q{Đã tạo N lệnh}. Lệnh mới
ở trạng thái Nháp. Hàng gấp và hạn giao chép từ đơn, ô Hạn hoàn thành sản xuất để trống chờ kế hoạch điền. Phần
mềm tự dựng khâu chuẩn bị giấy (cắt cuộn, tề giấy) theo nguồn giấy đang có trong kho.

Tab Lệnh SX có ba chip lọc Tất cả, Nháp, Sẵn sàng. Cột Hạn sản xuất tô màu theo số ngày còn dư so với hạn giao.

% ============================================================
\Needspace*{16\baselineskip}
\section{Rà một lệnh trước khi lập kế hoạch}\label{ra-lenh}

Bấm mã lệnh ở tab Lệnh SX để mở chi tiết. Việc ở chương này là làm cho lệnh đủ số để xưởng chạy được: đúng giấy,
đúng chuỗi công đoạn, mỗi bước có tổ và máy, bước cần khuôn có khuôn.

\rong{khsx-20-tong-quan.png}

\so{1} \oo{Còn thiếu N mục}: phần mềm tự đếm những gì lệnh còn thiếu. Bên dưới là từng mục kèm nút dẫn tới chỗ sửa.
\so{2} Ở đây là bước Khuôn hộp chưa chọn khuôn. Bấm tên bước, ngăn của bước đó mở ra.
\so{3} \oo{Sẵn sàng lập kế hoạch} mờ cho tới khi hết mục thiếu.
\so{4} Dãy ô số đầu trang: SL đặt, Vào máy, Giấy nguyên, Con/tờ in, số công đoạn, số món vật tư, hạn giao.
\so{5} \oo{Gia công trọn gói}: giao cả lệnh cho một nhà gia công (mục \ref{tron-goi}).
\so{6} \oo{Xoá lệnh}: chỉ có ở lệnh Nháp chưa giữ chỗ vật tư. Dòng đơn quay lại hàng chờ để lên lệnh lại.

Chip dưới mã lệnh dẫn về đơn, khách, PO, báo giá, phiếu tính giá. Có người khác vừa sửa lệnh thì nút \oo{Làm mới}
báo \q{Có thay đổi mới}, bấm mới nạp lại. Sửa ở tab nào mà chưa lưu thì đầu trang hiện \oo{Hoàn tác} và
\oo{Lưu thay đổi}.

Những mục \q{Còn thiếu} hay gặp: chưa có bài tính giá, thiếu loại giấy hay khổ nguyên của giấy, chưa có công đoạn,
thiếu ngày giao, có bước thuê ngoài chưa chọn nhà gia công, có công đoạn chưa gán tổ hoặc máy, có công đoạn cần
khuôn mà chưa chọn, chưa khai Con/tờ. Sách thì thêm Số trang, Trang mỗi tay, số mảnh xả.

\Needspace*{14\baselineskip}
\subsection{Thông tin chung}\label{thong-tin-chung}

\rong{khsx-21-thong-tin.png}

\so{1} Công tắc \oo{Hàng thường}, bật lên thành \q{Ưu tiên: Hàng GẤP}. Lệnh gấp đứng đầu khay Chờ xếp lịch.
\so{2} Hạn hoàn thành sản xuất: ngày xưởng phải làm xong, thường sớm hơn hạn giao khách vài ngày. Lịch dùng ngày này
để báo trễ, để trống thì lịch dùng hạn giao khách. Gõ ngày xong, bên dưới hiện thước đệm: \q{Đệm N ngày an toàn
trước hạn giao}, còn một ngày thì báo cận kề, trùng hoặc trễ hơn hạn giao thì báo đỏ.
\so{3} Ghi chú kế hoạch: ghi chú điều phối, dặn dò xưởng.

Phía dưới là khối Đơn hàng \& Nhân sự phụ trách và ô \oo{Lưu ý sản xuất (gửi xưởng)}, chép từ lưu ý của đơn.

\Needspace*{6\baselineskip}
\subsection{Quy cách}\label{quy-cach}

Tab Quy cách chỉ để xem: giấy, định lượng, khổ, cách in, số màu, số kẽm, con/tờ. Quy cách sai thì Kinh doanh sửa
phiếu tính giá, kế hoạch xoá lệnh Nháp rồi tạo lại từ hàng chờ.

\Needspace*{16\baselineskip}
\subsection{Công đoạn}\label{cong-doan}

Tab Công đoạn là chỗ kế hoạch sửa nhiều nhất. Phía trên là \oo{Nguồn giấy}: Tờ đúng khổ in, Tờ khổ nguyên, Cuộn
(gõ rộng cuộn rồi bấm Áp dụng) hoặc Mua mới. Đổi nguồn giấy thì phần mềm dựng lại bước cắt cuộn hay tề giấy cho
khớp, không bỏ tay các bước này được.

\rong{khsx-22-cong-doan.png}

Bảng có hai cách nhìn, \oo{Sơ đồ DAG} và \oo{Bảng danh sách}. Mỗi dòng là một bước: Loại (Máy, Tổ hay Thuê ngoài),
tổ và máy làm, số vào và ra kèm đơn vị, thời lượng máy chạy, bước đi trước (Tiền nhiệm). Số tròn đen cạnh tên là số
món vật tư khai ở bước. Nút ba chấm cuối dòng có Chuyển lên, Chuyển xuống, Chèn công đoạn sau, Bỏ bước.

Ảnh trên là lệnh 0007 sau khi đổi bước Cán màng bóng, mờ sang Thuê ngoài, giao Cơ sở Gia công Cán màng Thành Công.
Chuỗi khác bài tính giá thì phải ghi lý do ở ô cuối bảng, lý do vào Nhật ký của lệnh. Sửa xong bấm \oo{Lưu công
đoạn}. Dòng tổng cuối bảng cho biết tổng thời gian máy, dự kiến xong, còn mấy ngày tới hạn giao, có mấy bước thuê
ngoài.

Lệnh đang giữ chỗ vật tư thì tab này khoá, muốn sửa phải nhả hàng ở Kế hoạch vật tư (mục \ref{giu-hang}).

\Needspace*{14\baselineskip}
\subsection{Ngăn của một bước}\label{ngan-buoc}

\canh{0.47}{khsx-23-ngan-buoc.png}{%
Bấm tên bước, ngăn của bước mở bên phải, có năm tab: Cấu hình \& Số lượng, Phân công \& Thiết bị, Vật tư, Tiến độ
\& Thời gian, Phụ thuộc. Mũi tên trên đầu chuyển sang bước trước hay bước sau.

\so{1} Tab Phân công \& Thiết bị.

\so{2} Tổ phụ trách: chỉ có các tổ đã khai cho công đoạn này ở danh mục. Lúc ghi sản lượng, Việc của tổ chọn việc
khoán của tổ này.

\so{3} Máy sản xuất. Máy dùng chung nhiều tổ thì ghi thêm tên các tổ kia. Bước thuê ngoài thì chỗ này là ô chọn Nhà
gia công.

\so{4} Lời nhắn của sale về khuôn, chép từ phiếu tính giá.

\so{5} Chọn khuôn có sẵn của khách, hoặc \oo{+ Làm khuôn mới}. Chỉ hiện khuôn của đúng khách này.

\so{6} Sửa trong ngăn chưa lưu ngay. Bấm Xong để đóng ngăn rồi bấm Lưu công đoạn ở bảng chính.}

Danh sách tổ, máy và khuôn lấy từ danh mục, xem \kb{danh-muc-cong-doan}{Công đoạn},
\kbn{danh-muc-may}{Thiết bị \& Máy móc} và \kbn{danh-muc-khuon}{Khuôn}. Máy hay tổ nào không có trong ô chọn thì
nhờ người khai danh mục thêm.

\Needspace*{6\baselineskip}
\subsection{Vật tư, Tệp đính kèm, Nhật ký}\label{tab-khac}

Tab Vật tư gộp nhu cầu của cả lệnh theo nhóm: NVL chính, Phụ liệu, Dụng cụ, Khuôn, mỗi dòng ghi dùng ở bước nào.
Tab Tệp đính kèm giữ file gửi kèm lệnh. Tab Nhật ký ghi ai sửa gì, lúc nào, kèm lý do đổi chuỗi công đoạn.

\Needspace*{14\baselineskip}
\subsection{Sẵn sàng lập kế hoạch}\label{san-sang}

Hết mục thiếu thì nút Sẵn sàng lập kế hoạch sáng lên. Bấm là lệnh sang Sẵn sàng và vào khay Chờ xếp lịch.

\rong{khsx-24-san-sang.png}

\so{1} Chip trạng thái đổi sang Sẵn sàng.
\so{2} Chip \q{Chưa giữ chỗ vật tư}: lệnh chưa giữ hàng trong kho. Đây chỉ là nhắc, không chặn xếp lịch hay phát hành.
\so{3} \oo{Mở lại để sửa}: lệnh về Nháp. Lưu một thay đổi khi lệnh đang Sẵn sàng cũng đưa lệnh về Nháp.

\Needspace*{6\baselineskip}
\subsection{Hai băng nhắc danh mục đổi}\label{bang-nhac}

Lệnh lấy số của danh mục (định mức, thời lượng, vật tư từng bước) tại lúc tạo. Sau đó danh mục đổi thì đầu lệnh hiện
băng \q{Danh mục đã đổi sau lần lệnh này lấy số} kèm nút \oo{Cập nhật theo danh mục}. Không cập nhật thì lệnh vẫn
chạy với số cũ. Tương tự, tiêu chí KCS đổi sau lúc phát hành thì có băng nhắc kèm nút \oo{Cập nhật tiêu chí}.

% ============================================================
\Needspace*{16\baselineskip}
\section{Vật tư của lệnh}\label{vat-tu}

Mở \duong{Sản xuất > Kế hoạch vật tư}. Màn này gom vật tư của mọi lệnh chưa đóng: cần bao nhiêu, đã có bao nhiêu,
còn thiếu bao nhiêu, đang mua tới đâu.

\rong{khsx-30-vat-tu-truoc.png}

\so{1} Chip lọc theo tình trạng, số cạnh chip là số dòng.
\so{2} \oo{Xem theo} Lệnh (mỗi lệnh một nhóm) hoặc Mặt hàng (mỗi món một nhóm, gom mọi lệnh cần món đó). Ô tìm nhận mã
lệnh, tên món, mã phiếu mua.
\so{3} \oo{Giữ hàng} trên dòng nhóm của lệnh.
\so{4} \q{Cần mua}: kho không đủ, chưa có phiếu mua.
\so{5} \q{Chưa giữ}: kho có đủ mà lệnh chưa giữ.

Cột Đã có cộng ba phần: đã cấp cho xưởng, đang giữ trong kho, đang giữ ở hàng sắp về. Các tình trạng khác:
\q{Đang mua} (đã có phiếu mua chưa về), \q{Chờ hàng về} (đã giữ phần hàng sắp về), \q{Có trong kho} (đã giữ đủ), \q{Đã
cấp cho xưởng}, \q{Nên nhả hàng} (lệnh hết việc mà còn giữ hàng), \q{Chưa tính được} (thiếu số để tính nhu cầu).
\oo{Xuất Excel} lấy đúng các dòng đang lọc.

\Needspace*{14\baselineskip}
\subsection{Giữ hàng}\label{giu-hang}

Bấm \oo{Giữ hàng} để dành hàng trong kho cho lệnh, lệnh khác không lấy được nữa. Kho có bao nhiêu giữ bấy nhiêu, phần
thiếu tự giữ bù khi hàng về. Hai lệnh cùng chờ một món thì lệnh cần sớm hơn được giữ trước.

\rong{khsx-31-vat-tu-sau-giu.png}

\so{1} Món đã giữ đủ chuyển sang \q{Có trong kho}. Nút trên dòng nhóm đổi thành \oo{Nhả hàng}.

Lệnh đang giữ hàng thì khoá sửa số lượng, con/tờ, nguồn giấy, chuỗi công đoạn và không xoá được. Muốn sửa thì
\oo{Nhả hàng}, phần mềm hỏi lại một lần, nhả rồi không hoàn tác được. Giữ hàng không chặn xếp lịch hay phát hành:
lệnh chưa giữ vẫn phát hành được, chỉ mang chip cảnh báo.

\Needspace*{16\baselineskip}
\subsection{Đề nghị mua}\label{de-nghi-mua}

Dòng \q{Cần mua} thì tick ô đầu dòng, chọn được nhiều dòng của nhiều lệnh một lúc.

\rong{khsx-32-chon-de-nghi.png}

\so{1} Tick dòng cần mua. \so{2} Thanh dưới đáy đếm số dòng đã chọn, \oo{Bỏ chọn} để chọn lại. \so{3} Bấm \oo{Đề
nghị mua}.

\rong{khsx-33-form-ycmh.png}

Form \oo{Tạo yêu cầu mua hàng} mở ra, điền sẵn gần hết:

\so{1} Ngày cần hàng: bắt buộc. Gõ ngày hoặc bấm nút nhanh \so{2} Hôm nay, +3, +7, +14 ngày.
\so{3} Nội dung điền sẵn \q{Mua vật tư còn thiếu để chạy lệnh ...}, sửa được.
\so{4} Loại mua là \q{Cho lệnh SX}, không đổi được.
\so{5} Số lượng mua lấy bằng số còn thiếu, sửa được, kèm khổ của giấy.

Bấm Gửi yêu cầu. Thu mua nhận yêu cầu, chọn nhà cung cấp, lập đơn mua. Phần đó thuộc tài liệu Thu mua và Kho.

\rong{khsx-34-vat-tu-dang-mua.png}

\so{1} Cột Phiếu mua hiện mã yêu cầu, bấm để mở. \so{2} Bước của phiếu: Mới đề nghị, Đang lập đơn, Chờ duyệt, Đã đặt
hàng, Đã nhập kho, hoặc Đơn bị trả, chờ lập lại. \so{3} Dòng chuyển sang \q{Đang mua}. Hàng về kho thì lệnh tự giữ
bù và dòng sang \q{Có trong kho}.

Một món chỉ đề nghị mua được khi đã có nhà cung cấp khai bán món đó. Chưa có thì form báo \q{Vật tư chưa được nhà
cung cấp nào khai bán}, nhờ Thu mua khai món đó vào bảng giá của một nhà cung cấp rồi làm lại.

Bấm mã lệnh ở cột đầu mở ngăn của lệnh, có tab Vật tư, Phiếu mua, Lịch sử. Hai phiếu mua trùng nhau cho một món
thì ngăn hiện \q{Phiếu trùng nên huỷ} kèm nút \oo{Huỷ phiếu}, phải ghi lý do.

% ============================================================
\Needspace*{16\baselineskip}
\section{Xếp lịch và phát hành}\label{xep-lich}

Mở \duong{Sản xuất > Xếp lịch}. Lưới là các ngày, mỗi dòng một lệnh có giờ. Khay bên phải là lệnh Sẵn sàng chưa có
giờ.

\rong{khsx-40-khay.png}

\so{1} \oo{Chờ xếp lịch N}: bật tắt khay. \so{2} Cỡ nhìn Tuần, 2 tuần, Tháng. Mũi tên hai bên khoảng ngày để lùi tới,
\oo{Hôm nay} để về ngày hiện tại. \so{3} Thẻ lệnh trong khay: mã, tên, khách, hạn sản xuất, tổng giờ máy. Gấp và
hạn gần nhất đứng đầu. \so{4} Chip vật tư: chưa giữ, giữ một phần hay thiếu. \so{5} Lời nhắc cách kéo thả.
\so{6} \oo{Khoá máy}: chặn một khoảng thời gian của máy (mục \ref{khoa-may}).

Ô tìm (phím tắt \oo{/}) nhận mã, tên, khách. \oo{Lọc} theo trạng thái (Chưa phát hành, Đã phát hành, Đã xong),
khách, hàng gấp. \oo{Hiển thị} chọn thông tin hiện trên thanh.

\Needspace*{14\baselineskip}
\subsection{Kéo lệnh vào lịch}\label{keo-tha}

Kéo thẻ lệnh từ khay thả vào cột ngày muốn bắt đầu.

\rong{khsx-41-dang-tha.png}

Thả vào ngày nghỉ hoặc ngoài ca thì lệnh tự sang ca kế tiếp. Lệnh chạy từ đầu ca (ở đây Ca 2, 14:00), các bước nối
nhau theo thời lượng máy của từng bước.

\rong{khsx-42-sau-tha.png}

Dòng lệnh vừa thả hiện trên lưới, khay còn 0. \so{2} Thanh lệnh bắt đầu 14:00 thứ Sáu 6/11. Băng tóm tắt trên lưới
đếm tổng giờ máy trong khoảng đang xem và có hai chip lọc nhanh: \so{3} \q{Có thể xong muộn hơn} (có bước chưa tính
được giờ) và \so{4} \q{Chưa phát hành}. Lệnh trễ thì có thêm chip \q{Trễ hạn}, thanh có đuôi sọc đỏ và chip \q{trễ N
ngày}.

Đặt giờ không đổi trạng thái lệnh, chưa có gì xuống xưởng. Dời lệnh có bốn cách:

\begin{itemize}
\item Kéo thanh sang ngày khác (dời nguyên ngày).
\item Gõ giờ ở ô Bắt đầu trong ngăn lệnh.
\item Bấm \oo{← 15} hoặc \oo{15 →} trong ngăn để lùi tới 15 phút.
\item Chọn thanh rồi bấm phím mũi tên trái phải (15 phút), giữ Shift để nhảy một ca.
\end{itemize}

Giờ rơi ra ngoài ca thì tự trượt sang ca kế. Dời nhầm thì \oo{Ctrl+Z} hoặc nút Hoàn tác trên thông báo, lùi được 20
bước. \oo{Bỏ lịch} trong ngăn đưa lệnh về khay.

\Needspace*{16\baselineskip}
\subsection{Ngăn lệnh trên lịch}\label{ngan-lich}

Bấm vào thanh hoặc mã lệnh, ngăn lệnh mở ra.

\rong[0.8\linewidth]{khsx-43-ngan-lenh.png}

\so{1} Ô Bắt đầu: ngày và giờ, kèm nút lùi tới 15 phút. \so{2} Bước chưa tính được giờ (thường là bước tổ làm tay
chưa khai thời lượng) thì ngày xong có thể muộn hơn số đang hiện. Dãy ô bên cạnh: Sớm nhất xong, Từ đầu tới xong,
Máy chạy, Hạn xong sản xuất, Hạn giao khách, kèm số ngày còn dư. \so{4} Bước thuê ngoài ghi \q{Gia công ngoài}
thay cho giờ. \so{5} Chip vật tư của lệnh. \so{6} \oo{Bỏ lịch}. \so{7} \oo{Phát hành}.

Nút \oo{Mở ở Kế hoạch sản xuất} trên đầu ngăn mở chi tiết lệnh. Lệnh đã phát hành thì nút này đổi thành \oo{Mở hồ
sơ lệnh}.

\Needspace*{12\baselineskip}
\subsection{Phát hành}\label{phat-hanh}

Bấm \oo{Phát hành}. Phần mềm không hỏi lại. Lệnh cùng nhóm thành phẩm hoặc phụ thuộc chéo nhau được phát hành cùng
nhau, khi đó nút ghi \q{Phát hành cả cụm, N lệnh}.

\cat[0.8\linewidth]{0 0 0 1150}{khsx-44-toast-phat-hanh.png}

Phát hành xong:

\begin{itemize}
\item Lệnh sang Đã phát hành. Tổ thấy việc ở Việc của tổ theo phân công trong lệnh.
\item Mỗi bước có vật tư tự lập một yêu cầu xuất kho, thủ kho nhận thông báo. Thông báo ghi \q{Đã gửi kho N yêu cầu
xuất}. Món nào chưa tự lập được thì thông báo nói rõ, tổ xin tay ở Việc của tổ.
\item Bước thuê ngoài sang \q{Chờ mang đi} (chương \ref{gia-cong}).
\end{itemize}

Thiếu vật tư hay chưa giữ hàng không chặn phát hành. Người kế hoạch tự quyết dựa trên chip cảnh báo.

\Needspace*{30\baselineskip}
\subsection{Sau khi phát hành: Thu hồi, Phát hành cập nhật}\label{thu-hoi}

\rong[0.8\linewidth]{khsx-45-ngan-da-phat-hanh.png}

\so{1} \oo{Thu hồi phát hành}: rút lệnh khỏi xưởng, lệnh về \q{Có lịch, chưa phát hành}, yêu cầu xuất kho chưa ghi
sổ bị huỷ. Chỉ thu hồi được khi chưa việc nào của lệnh bắt đầu. Phải ghi lý do, ít nhất ba ký tự.

\so{2} \oo{Phát hành cập nhật}: dời giờ hay đổi máy sau khi đã phát hành mà không rút lệnh về. Chỉ áp cho việc chưa
bắt đầu, phải ghi lý do, tổ nhận được và xác nhận lại phân công.

\cat[0.5\linewidth]{190 410 703 410}{khsx-47-thu-hoi.png}

Tab \oo{Lịch sử phát hành} ghi từng phiên bản, lần thu hồi kèm lý do. Bấm chọn hai phiên bản để so giờ và máy từng
bước.

\cat[0.8\linewidth]{0 785 0 0}{khsx-46-lich-su-phat-hanh.png}

Ô quyền \q{Phát hành lịch} ở Xếp lịch quyết ai được Phát hành, Phát hành cập nhật, Thu hồi. Người chỉ có Thao tác
thì kéo thả, dời, bỏ lịch được, không phát hành được.

\Needspace*{8\baselineskip}
\subsection{Khoá máy}\label{khoa-may}

Máy bảo trì, hỏng hay nghỉ thì bấm \oo{Khoá máy} trên đầu màn. Mỗi khoảng khoá gồm Máy, Từ, Đến, Lý do (Bảo trì,
Hỏng, Nghỉ, Khác), Ghi chú. \oo{Thêm khoảng khoá} để khai nhiều khoảng. Lệnh đã xếp đè lên khoảng khoá không tự dời,
lịch chỉ cảnh báo để kế hoạch dời tay.

% ============================================================
\Needspace*{16\baselineskip}
\section{Gia công ngoài}\label{gia-cong}

Có hai cách đưa hàng ra ngoài. Thuê ngoài một phần: một hay vài bước của lệnh do nhà gia công làm, các bước khác
vẫn chạy ở xưởng. Gia công trọn gói: cả lệnh giao cho một nhà gia công. Nhà gia công là nhà cung cấp có bật
\q{Nhận gia công}, do Thu mua khai.

Người kế hoạch cũng là người mang hàng đi và ghi số nhận về. Mọi việc đó làm ở khối \oo{Gia công ngoài} trên chi tiết
lệnh, ở màn Kế hoạch sản xuất.

\Needspace*{14\baselineskip}
\subsection{Thuê ngoài một phần}\label{mot-phan}

Ở tab Công đoạn, đổi Loại của bước sang Thuê ngoài và chọn nhà gia công trong ngăn bước (mục \ref{ngan-buoc}). Phát
hành xong, các bước thuê ngoài liền nhau cùng một nhà gia công thành một lần gia công.

\rong{khsx-50-gia-cong-mot-phan.png}

\so{1} Chip tình trạng: Chờ mang đi, Đang ở nhà gia công, Đã nhận về, Đã huỷ. Dải chấm bên cạnh là ba nấc đó.
\so{2} Nhà gia công kèm số điện thoại. \so{3} Bước trước chưa bàn giao hàng sang thì chưa mang đi được.
\so{4} \oo{Lịch sử}: ai làm gì, lúc nào.

Mang hàng đi:

\begin{itemize}
\item Bước thuê ngoài là bước đầu tiên: gõ Số mang đi rồi bấm \oo{Đã mang đi}.
\item Có bước trước: số mang đi lấy theo số tổ trước đã bàn giao. Tổ bàn giao thêm thì bấm \oo{Mang thêm}.
\end{itemize}

Hàng về thì bấm \oo{Nhận hàng về}, gõ Số nhận về, chọn \oo{Hàng về đâu} và bước nhận hàng, rồi \oo{Chốt số nhận về}.
Còn bước sau thì hàng \q{Về xưởng làm tiếp}, tổ của bước sau thấy hàng để nhận. Bước thuê ngoài là bước cuối thì hàng
nhập kho thành phẩm.

Chốt nhầm thì \oo{Mở lại} khi lần gia công đang Đã nhận về. Mở lại bị khoá khi đã có phiếu chi cho lần gia công, tổ
nhận đã xác nhận, kho đã lập phiếu, hoặc lệnh đã đóng.

KCS không kiểm bước thuê ngoài, số chốt nhận về là số chính thức. Tiền gia công do Kế toán chi: dòng Chi tiền ghi
\q{Chờ kế toán chi}, chi xong thì hiện mã phiếu chi.

\Needspace*{16\baselineskip}
\subsection{Gia công trọn gói}\label{tron-goi}

Lệnh chưa phát hành, không nằm trong bài ghép, thì có nút \oo{Gia công trọn gói} trên đầu chi tiết lệnh.

\canh{0.42}{khsx-51-dat-tron-goi.png}{%
\so{1} Nhà gia công.

\so{2} Số đặt, mặc định bằng số lượng lệnh.

\so{3} \oo{Xưởng cấp giấy}: tick khi công ty lo giấy theo phiếu tính giá. Bỏ tick khi nhà gia công tự lo giấy.

\so{4} Bấm \oo{Đặt gia công trọn gói}.

Lệnh trọn gói không xuống tổ, không vào Xếp lịch. Lệnh sang Đã phát hành ngay, chuỗi công đoạn chỉ còn một bước
\q{Gia công trọn gói} ở nhà gia công đã chọn.}

\rong{khsx-52-tron-goi-sau-dat.png}

\so{1} Lệnh đã phát hành. \so{2} Lần gia công sang thẳng \q{Đang ở nhà gia công}, không qua nấc Chờ mang đi.
\so{3} Xưởng cấp giấy mà chưa chọn giấy thì có dải \q{Chưa cấp giấy cho nhà gia công} kèm nút \oo{Chọn giấy}.
\so{4} \oo{Huỷ gia công}: lệnh về Nháp, phải ghi lý do. Kho đã lập phiếu xuất giấy cho nhà gia công thì không huỷ được.

\Needspace*{14\baselineskip}
\subsection{Cấp giấy cho nhà gia công}\label{cap-giay}

Bấm \oo{Chọn giấy}.

\rong{khsx-53-chon-giay.png}

\so{1} Các khổ của loại giấy này đang có trong kho, kèm số tồn. Khổ đúng lệnh được chọn sẵn.
\so{2} Số tờ xin xuất, điền sẵn theo lệnh. \so{3} Kho thiếu thì phần mềm báo nhưng vẫn cho gửi, kho xuất trước phần
đang có. \so{4} \oo{Gửi đề nghị xuất giấy}. \oo{Để sau} thì đóng lại, dải nhắc vẫn còn.

\rong{khsx-54-da-gui-de-nghi-giay.png}

\so{1} Tình trạng phía kho: Chờ kho soạn, Kho đã nhận, Kho đang soạn, Đã xuất một phần, Đã xuất đủ. Bấm mã đề nghị
(DNX...) để xem. \so{2} \oo{Nhận hàng về} mờ tới khi kho ghi sổ phiếu xuất giấy, dòng chữ bên dưới nói rõ đang chờ gì.

Hàng trọn gói về thì bấm Nhận hàng về, gõ số, hàng nhập thẳng kho thành phẩm.

% ============================================================
\Needspace*{16\baselineskip}
\section{Theo dõi}\label{theo-doi}

Kinh doanh và xưởng hỏi liên tục \q{lệnh này tới đâu}. Có hai màn để trả lời: Theo dõi sản xuất nhìn theo đơn, Hồ sơ
lệnh sản xuất nhìn theo lệnh. Phía Kinh doanh đọc cùng màn Theo dõi sản xuất, xem \kd{theo-doi}{Trả lời khách: hàng
tới đâu rồi}.

\Needspace*{14\baselineskip}
\subsection{Theo dõi sản xuất}\label{theo-doi-don}

Mở \duong{Sản xuất > Theo dõi sản xuất}. Mỗi dòng là một đơn.

\rong{khsx-60-theo-doi.png}

\so{1} Chip Tất cả, Cần để ý, Đang làm, Chờ giao, Đã giao đủ. \q{Cần để ý} gom đơn đang vướng, đã trễ hoặc còn ba ngày
tới hạn mà chưa đủ hàng. \so{2} Bốn chặng Lệnh, SX, Kho, Giao. \so{3} Chấm đỏ là chặng mặt hàng chậm nhất của đơn đang
đứng: DH003 còn mặt hàng chưa xuống xưởng nên đứng ở Lệnh. \so{4} Hạn giao kèm \q{còn N ngày} hoặc \q{trễ N ngày}.
\so{5} Đang vướng: Trễ hạn giao, Sự cố máy, Tạm dừng, Thiếu vật tư, KCS không đạt, Đóng lệnh thiếu, Chưa có lệnh.

Ô tìm (Ctrl+K) nhận số đơn, khách, mặt hàng, mã lệnh. Kỳ lọc theo ngày chốt đơn hoặc hạn giao khách. \oo{Lọc} theo
khách, người phụ trách.

Bấm một đơn để mở ngăn. Tab Tổng quan liệt kê mặt hàng với số đặt, nhập kho, có hàng, đã giao. Bấm mũi tên trước
mặt hàng để bung các lệnh của nó, mỗi lệnh có dải bước và chỗ đang vướng. Cột phải tóm tắt hạn giao, ngày giao xưởng,
danh sách vướng, tình trạng giao của từng mặt hàng.

\rong[0.8\linewidth]{khsx-61-ngan-don.png}

\Needspace*{14\baselineskip}
\subsection{Hồ sơ lệnh sản xuất}\label{ho-so-lenh}

Mở \duong{Sản xuất > Hồ sơ lệnh sản xuất}. Chỉ có lệnh đã phát hành hoặc đã đóng.

\rong{khsx-62-ho-so-lenh.png}

\so{1} Chip Tất cả, Đang sản xuất, Sau sản xuất, Đã giao đủ. \so{2} Các bước: chấm xanh là bước đã xong, chấm cam là
bước đang làm. \so{3} Trạng thái theo khâu: Đang sản xuất, Đang KCS, Chờ nhập kho, Sẵn sàng giao, Đã giao đủ.
\so{4} Lệnh trọn gói chỉ có một dòng \q{Gia công trọn gói ở ...}. \oo{Lọc} có mục Gia công ngoài để chỉ xem lệnh có
bước thuê ngoài.

Bấm một lệnh để mở hồ sơ.

\rong[0.8\linewidth]{khsx-63-ngan-ho-so.png}

\so{1} Bốn tab Quy cách, Sản lượng, Kho và giao, Lịch sử. \so{2} \oo{In phiếu công nghệ} cho xưởng.
\so{3} Cột phải \q{Hàng đang tới đâu}: dải bước dọc, bước thuê ngoài ghi tên nhà gia công, cuối dải là nhập kho thành
phẩm. \so{4} Dải vàng trên cùng là chỗ đang vướng của lệnh.

Tab Sản lượng cho số tổng (Đặt, Làm ra, Đạt KCS cuối, Nhập kho, Đã giao) và bảng Từng bước: làm ra trên kế hoạch,
xong lúc nào, chưa xong thì giờ hẹn theo lịch.

\rong[0.8\linewidth]{khsx-64-san-luong.png}

Tab Kho và giao cho biết thành phẩm đã vào kho bao nhiêu, đã xuất cho khách bao nhiêu. Tab Lịch sử ghi mọi việc của
lệnh từ lúc tạo.

% ============================================================
\Needspace*{10\baselineskip}
\section{Khi lệnh xong}\label{lenh-xong}

Kế hoạch không có nút đóng lệnh. Bước cuối làm xong thì KCS kiểm, gửi thành phẩm vào kho, rồi đóng lệnh theo nhóm
thành phẩm. Thành phẩm chỉ vào kho khi có người giữ quyền kiểm KCS, xem \kb{danh-muc-kcs-quyen}{Ai được kiểm KCS}.

Làm ra ít hơn số đặt thì KCS chỉ đóng được khi có quyền \q{Đóng lệnh thiếu}. Lệnh đóng thiếu thì mặt hàng hiện
\q{Đóng lệnh thiếu} ở cột Đang vướng của Theo dõi sản xuất, đơn giao phần đã có. Phần mềm không tự tạo lệnh làm bù.

Hàng gia công trọn gói và bước thuê ngoài cuối chuỗi không qua KCS: số chốt nhận về là số nhập kho.

% ============================================================
\Needspace*{10\baselineskip}
\section{Câu hỏi hay gặp}\label{hoi-dap}

\begin{description}[style=nextline,leftmargin=0pt,font=\normalfont\tieude\small]
\item[Kinh doanh nói đã chốt đơn mà Hàng chờ tiếp nhận không có?]
Đơn cần cọc chưa đủ cọc thì chưa sang sản xuất.
Hỏi Kế toán, hoặc xem đơn ở \kd{don-coc}{Chờ cọc}. Đơn đã có lệnh cho mọi dòng cũng rời hàng chờ.

\item[Nút Sẵn sàng lập kế hoạch mờ?]
Đọc \oo{Còn thiếu N mục} ở đầu chi tiết lệnh, bấm từng mục để tới chỗ sửa
(mục \ref{ra-lenh}).

\item[Không sửa được số lượng hay chuỗi công đoạn?]
Lệnh đang giữ hàng. Nhả hàng ở Kế hoạch vật tư (mục
\ref{giu-hang}). Lệnh Sẵn sàng thì bấm Mở lại để sửa trước. Lệnh đã phát hành thì chuỗi công đoạn khoá hẳn, chỉ
dời giờ hay đổi máy bằng Phát hành cập nhật.

\item[Lệnh không có trong khay Chờ xếp lịch?]
Lệnh còn Nháp, đã có giờ trên lưới, hoặc là lệnh trọn gói (trọn gói
không xếp lịch).

\item[Phát hành khi chưa đủ vật tư có được không?]
Được. Phần mềm chỉ cảnh báo bằng chip. Món chưa có thì kho chưa
xuất được, tổ chờ hoặc xin tay.

\item[Phát hành rồi mới thấy sai?]
Chưa việc nào bắt đầu thì Thu hồi phát hành, sửa, xếp lại, phát hành lại. Đã có
việc bắt đầu thì chỉ còn Phát hành cập nhật cho các việc chưa bắt đầu.

\item[Thanh lệnh có đuôi sọc đỏ?]
Lệnh xong muộn hơn Hạn hoàn thành sản xuất (hoặc hạn giao khách nếu lệnh không
có hạn sản xuất). Dời sớm lên, đổi máy, hoặc báo Kinh doanh.

\item[Đề nghị mua báo vật tư chưa được nhà cung cấp nào khai bán?]
Nhờ Thu mua khai món đó vào bảng giá của một nhà
cung cấp, rồi đề nghị lại.

\item[Nút Nhận hàng về ở lệnh trọn gói mờ?]
Xưởng cấp giấy mà kho chưa xuất giấy cho nhà gia công. Xem tình trạng
kho ở khối Gia công ngoài (mục \ref{cap-giay}).

\item[Ai đóng lệnh?]
KCS. Kế hoạch chỉ theo dõi tới lúc lệnh sang Đã đóng (chương \ref{lenh-xong}).
\end{description}

\end{document}
